\section{Oracle-Grounded Training: Quasi-Sentence Supervision on the Manifesto Corpus}\label{sec:manifesto-qsentence}

The pipeline of Section~\ref{sec:manifesto-llm} is prompt-only: the
oracle enters only at the root, through expert-survey means, and the
local laws appear only as audit probes. The Manifesto
corpus supports a stronger design. Trained human coders assign every
quasi-sentence (a clause-level unit; roughly one statement per unit) a
category code from the Comparative Manifesto Project (CMP) coding
scheme. Because those codes aggregate additively over spans, they
define a gold oracle value at \emph{every node of every tree}: each
leaf, each merge, and the root. This is the setting the controlled
sections train in, now on real political text: $g$ and $f$ can be
supervised at every node, and the local laws become training signals
rather than audit probes. Below are the two label systems,
the \texttt{treepo} implementation that runs the experiments, and the
results, including the leaf-size ladder (Appendix~\ref{app:qsentence-leaf-ladder})
and gold-seeded summarization (Appendix~\ref{app:gold-seeding}).

\subsection{Two Label Systems: Documents and Quasi-Sentences}\label{sec:qsentence-labels}

The Manifesto measurement problem comes with two distinct layers of
ground truth, and the distinction does real work in what follows.

\paragraph{Document-level labels.} Each platform carries two kinds of
document-level target. The first is the pair used in
Section~\ref{sec:manifesto-llm}: the \citet{BenoitEtAl2025}
expert-survey means on six seven-point policy scales, available for the
$235$-platform benchmark. The second is the Manifesto Project's own
published scores \citep{ManifestoProject2025a}: the RILE left--right
index on $[-100, 100]$ and per-domain salience shares, available for
every coded platform. RILE is a deterministic function of the CMP
category counts: with $R$ the count of quasi-sentences bearing one of
$13$ designated right codes, $L$ the count bearing one of $13$ left
codes \citep{LaverBudge1992}, and $N$ the total count of coded
non-header quasi-sentences,
\begin{equation}\label{eq:rile}
\mathrm{RILE} \;=\; 100\cdot\frac{R - L}{N},
\qquad
\mathrm{RILE}_{\mathrm{norm}} \;=\; \frac{\mathrm{RILE} + 100}{200}\in[0,1].
\end{equation}
In words: RILE is the signed share of coded statements that lean right,
so it inherits the additivity of the counts it is built from. The seven
domain shares $\mathrm{dom}_1,\dots,\mathrm{dom}_7$ are the analogous
count ratios for the seven CMP handbook policy domains (external
relations, freedom and democracy, political system, economy, welfare
and quality of life, fabric of society, social groups).

\paragraph{Quasi-sentence gold labels.} The Manifesto Corpus release
provides the underlying per-quasi-sentence annotations for $2{,}157$
platforms, about $2.27$M coded spans: each quasi-sentence carries one
CMP category code from a $56$-category scheme, assigned by a trained
human coder under the project's handbook. These are the gold labels.
For any contiguous span $x$ of a platform, aggregating the codes of the
quasi-sentences inside $x$ and applying Equation~\eqref{eq:rile} yields
a gold oracle value $\fstar(x)$: an $8$-vector
$(\mathrm{RILE}_{\mathrm{norm}}, \mathrm{dom}_1, \dots, \mathrm{dom}_7)$
we call the \emph{compact state}. Because the aggregation is a count
sum, the oracle is exactly mergeable in the sense of
Section~\ref{sec:mergeable-sketches}: the compact state of a parent
span is a known function of its children's counts. Every node of every
C-Tree over a coded platform therefore has a gold target, which is the
structural property the controlled Markov and HLL sections relied on
and that generic long-document corpora lack.

\paragraph{Rollup validity.} The first thing to check is that the
bottom of this label system reproduces the top. Rolling the
quasi-sentence codes up to the document level and correlating with the
published MPDS document scores across the $2{,}157$ coded platforms
gives Pearson $r = 0.9975$ for RILE and $r \in [0.924, 0.999]$ across
the seven domain shares.\footnote{The residual is bookkeeping, not
noise: the gold rollup uses the non-header quasi-sentence count as its
denominator and folds handbook-v5 four-digit subcodes into their
three-digit parents, while the published per-category percentages use
MPDS's own denominator conventions. Correlations computed by
\texttt{scripts/check\_qsentence\_doc\_label\_correlation.py}; full
per-category table in the artifact
\texttt{outputs/qsentence\_doc\_label\_correlation/}.} The
document-level published label is, to numerical precision, the additive
rollup of the quasi-sentence labels. Local-to-global reconstruction is
therefore a well-posed target on this corpus: a method that recovers
the right local states and merges them lawfully recovers the published
document score.

\paragraph{What the codes cannot reach.} The same check run against
the \citet{BenoitEtAl2025} expert-survey means gives a much lower
ceiling: ordinary least squares from the eight compact-state features
to the expert mean attains five-fold cross-validated $R^2$ between
$0.03$ (decentralization) and $0.47$ (immigration). The expert-survey
construct is a holistic judgment that the additive category counts do
not determine. This splits the measurement problem into two targets
with different supervision economics. The \emph{internal} target (the
teacher rollup of gold quasi-sentence codes) is fully decomposable and
supervisable at every node. The \emph{external} target (expert-survey
means) requires the model to read text, and quasi-sentence codes alone
cannot reach it. The experiments below report both, following the
principle that a low error against one target with no correlation
against the other is a measurement failure, not a success.

\subsection{Bundles, Splits, and the \texttt{treepo} Harness}\label{sec:qsentence-harness}

The experiments run on frozen \emph{labeled-tree bundles}. A bundle
fixes a leaf granularity $\ell \in \{1, 2, 4, 8, 16\}$ quasi-sentences
per leaf, partitions each platform into contiguous leaves of $\ell$
quasi-sentences, builds the binary merge tree, and stamps every node
with its gold compact state via the count aggregation above; no LLM is
involved in labeling. The corpus is $218$ Western European platforms
under a fixed split ($140$ train, $30$ validation, $48$ test) shared
across all leaf sizes, so a leaf-size comparison never mixes documents.
Granularity changes the supervision surface by an order of magnitude:
the $\ell = 1$ bundle has about $437{,}000$ labeled nodes, the
$\ell = 16$ bundle about $28{,}000$.

The implementation is the \texttt{treepo} package. Its public entry
point is \texttt{treepo.fit}, which takes a declarative spec (family,
data, schedule) and runs the alternating $f/g$ ladder of
Section~\ref{sec:manifesto-llm} for any \emph{family}: an object
implementing the four-method contract
\begin{lstlisting}[language=Python]
class FamilyRuntime(Protocol):
    name: str
    def train_f(self, *, f_init, g, traces, output_dir, iteration): ...
    def train_g(self, *, g_init, f, traces, output_dir, iteration): ...
    def score_roots_with_f(self, *, f, g, trees): ...
    def validate_artifact(self, *, kind, artifact): ...
\end{lstlisting}
The contract is the operator interface of
Section~\ref{sec:framework} made executable: a family supplies its own
parameterization of the summarizer $g$ and readout $f$, and the
surrounding loop (alternation schedule, node sampling, metrics, audit
logging) is family-agnostic. Two families instantiate it here, chosen
to sit at opposite ends of the parameterization spectrum:
\begin{itemize}[leftmargin=1.5em]
    \item \textbf{LLM prompt-program family} (\texttt{dspy}): $f$ and
          $g$ are prompt programs over a fixed open-weight checkpoint,
          optimized by prompt search against node-level rewards. The
          state passed between nodes is text.
    \item \textbf{Neural-operator family} (\texttt{fno}): leaves are
          embedded once by a small embedding model
          (\textsc{embeddinggemma-300m}), and $f$ and $g$ are Fourier
          Neural Operators trained by gradient descent on node-level
          losses, exactly as in Section~\ref{sec:fno-primer}. The state
          is a learned vector.
\end{itemize}
Both families consume the same bundles and the same splits, so their
numbers are directly comparable.

A minimal run, in the style of a package vignette, is three calls:
build the trees, declare the backend, fit.
\begin{lstlisting}[language=Python]
from treepo import fit
from treepo.tasks.manifesto import (
    make_manifesto_replication_trees, manifesto_prompt_template)

train = make_manifesto_replication_trees(
    split="train", leaf_unit_count=8, doc_unit_kind="qsentence")
evaltrees = make_manifesto_replication_trees(
    split="test", leaf_unit_count=8, doc_unit_kind="qsentence")

result = fit({
    "family": "dspy",              # or "fno"
    "train_data": train,
    "eval_data": evaltrees,
    "backend_config": {
        "model": "openai/nvidia/Gemma-4-31B-IT-NVFP4",
        "prompt_template": manifesto_prompt_template(),
    },
    "axis": {"axis_kind": "leaf_unit_count", "axis_value": 8,
             "max_iterations": 2},
})
print(result.status, result.metrics["internal_f_mae"])
\end{lstlisting}
The research sweeps behind the tables below drive the same machinery
through a CLI that exposes the leaf-size axis and the family switch;
the two substrate legs of the ladder differ in a single flag:
\begin{lstlisting}[language=bash]
# LLM prompt-program leg
python scripts/run_manifesto_qsentence_dspy_ladder.py \
  --family dspy --leaf-qsentences "1,2,4,8,16" \
  --max-iterations 2 --target-dimensions all \
  --dspy-optimizer gepa \
  --dspy-model "openai/nvidia/Gemma-4-31B-IT-NVFP4"

# neural-operator leg: same bundles, same splits
python scripts/run_manifesto_qsentence_dspy_ladder.py \
  --family fno --leaf-qsentences "1,2,4,8,16" \
  --embedding-model google/embeddinggemma-300m \
  --fno-target-dimension rile --fno-epochs 8
\end{lstlisting}
Each run leaves behind the same artifact schema (per-stage grid
summaries, per-node prediction records, iteration histories), which is
what makes the substrate comparisons in
Appendix~\ref{app:qsentence-leaf-ladder} a table join rather than a
reconciliation exercise.

Metrics follow the two-target split of
Section~\ref{sec:qsentence-labels}. \emph{Internal} Pearson and MAE
compare root predictions to the teacher rollup (the gold compact state
of the root); \emph{external} Pearson compares them to the published
MPDS document score. All correlations below are per-dimension:
pooling the eight dimensions into one correlation inflates the number
by between-dimension mean separation (the pooled statistic rewards a
model for predicting that economy is discussed more than
decentralization, which is not measurement), so a per-dimension
report is the honest one.

\subsection{Results: The Learned Merge Composes Where the Signal Is Dense}\label{sec:qsentence-results}

\paragraph{Leaf-size ladder.} Table~\ref{tab:qsentence-fno-ladder}
(abridged; full grid in Appendix~\ref{app:qsentence-leaf-ladder})
reports the neural-operator family on RILE across leaf sizes, test
split, $n = 48$. Two stages matter: the $f$-stage (readout trained on
leaf states, merge untrained) and the $g$-stage (learned merge trained
against the current $f$).

\begin{table}[htbp]
    \centering
    \begin{tabular}{rccc}
    \toprule
    leaf ($\ell$ qsent.) & $f$-stage ext.\ $r$ & $g$-stage ext.\ $r$ & $g$-stage MAE \\
    \midrule
    1  & $0.450$ & $0.479$ & $0.069$ \\
    2  & $0.553$ & $0.537$ & $0.085$ \\
    4  & $0.537$ & $0.538$ & $0.066$ \\
    8  & $0.664$ & $\mathbf{0.718}$ & $0.058$ \\
    16 & $0.645$ & $0.672$ & $0.070$ \\
    \bottomrule
    \end{tabular}
    \caption{Neural-operator (FNO) family on RILE across quasi-sentence
    leaf sizes. External Pearson $r$ against the published MPDS RILE
    (normalized scale), test split ($n=48$), $g$-stage $= f^2g^2$.
    Appendix~\ref{app:qsentence-leaf-ladder} gives the full grid with
    internal metrics and the LLM-family counterpart.}
    \label{tab:qsentence-fno-ladder}
\end{table}

Two features of the table carry the message. First, the learned merge
helps at the granularities where it has signal to compose: at
$\ell = 8$ the $g$-stage improves on the $f$-stage from $0.664$ to
$0.718$, and the peak of the ladder sits there rather than at the
coarsest leaf. Second, the profile over leaf size is a tension the
fixed-prompt pipeline of Section~\ref{sec:manifesto-llm} did not have:
smaller leaves mean more supervised nodes but longer merge chains for
a \emph{trained} operator to keep lawful, and at a fixed training
budget the product peaks at an intermediate granularity. The
prompt-only pipeline's chunk-invariance ($\pm 0.027$ macro over a
$16\times$ sweep) is the behavior of a fixed, already-competent
summarizer; the trained substrate's leaf-size profile is the cost of
learning that summarizer from local supervision alone.

\paragraph{Per-dimension substrate comparison.} Running both families
on the same $\ell = 8$ bundle and scoring each of the eight compact
dimensions separately gives Table~\ref{tab:qsentence-perdim} (internal
Pearson, $g$-stage). The comparison needs one confound named before
the numbers can be read. The LLM family's $f$-stage sees the gold leaf
label in its input trace, and its readout learns to echo it
(internal $r = 1.000$ exactly at the $f$-stage); the FNO family
predicts from text embeddings alone. LLM cells are therefore
upper bounds inflated by label echo, and FNO cells are honest
predictions.

\begin{table}[htbp]
    \centering
    \begin{tabular}{lccccc}
    \toprule
    dimension & teacher mean & FNO $f$ & FNO $g$ & dgemma $g$ & Gemma-4 $g$ \\
    \midrule
    RILE       & $0.47$  & $0.67$ & $\mathbf{0.72}$ & $0.12$  & $0.25$ \\
    domain\_5  & $0.30$  & $0.67$ & $\mathbf{0.72}$ & $-0.04$ & $0.37$ \\
    domain\_4  & $0.20$  & $0.64$ & $\mathbf{0.62}$ & $0.41$  & $0.04$ \\
    domain\_6  & $0.15$  & $0.63$ & $\mathbf{0.63}$ & $-0.27$ & $-0.11$ \\
    \midrule
    domain\_3  & $0.125$ & $0.02$ & $0.02$  & $0.17$  & $0.32$ \\
    domain\_7  & $0.09$  & $-0.13$ & $-0.13$ & $-0.13$ & $0.18$ \\
    domain\_1  & $0.07$  & $0.00$ & $0.00$  & $-0.13$ & $-0.31$ \\
    domain\_2  & $0.07$  & $0.00$ & $0.00$  & $-0.06$ & $0.13$ \\
    \bottomrule
    \end{tabular}
    \caption{Per-dimension internal Pearson at the $g$-stage, $\ell=8$
    bundle, test split. ``Teacher mean'' is the mean gold share of the
    dimension (a signal-density proxy). The horizontal rule separates
    dense dimensions (mean $\gtrsim 0.13$) from sparse ones. LLM
    columns (dgemma = DiffusionGemma runtime, Gemma-4 = Gemma-4-31B)
    are inflated by gold-label echo at the leaf stage; see text.}
    \label{tab:qsentence-perdim}
\end{table}

The table is a signal-density story rather than a substrate horse race.
On every dense dimension (gold share $\gtrsim 0.13$: RILE,
domains~4--6) the neural-operator family both learns the leaf readout
($f$ between $0.63$ and $0.67$) and composes it through the learned
merge ($g$ between $0.62$ and $0.72$), beating both LLM merges. On the
sparse dimensions ($87$--$93\%$ of leaves carry a zero share) every
substrate fails: the FNO readout collapses to predicting zero, and the
apparent LLM wins there are the label echo just described, not reading.
The LLM merge's failure on dense dimensions has a specific, fixable
cause: the prompt search optimizes a mean-over-dimensions accuracy
reward, and with most dimensions sparse, predicting each dimension's
corpus mean already scores about $0.90$ on that reward with roughly
zero within-dimension correlation. The optimizer finds that degenerate
optimum. This is a concrete instance of the projection story of
Section~\ref{sec:theorems}: an objective that does not price the local
laws lets the learned operator drift to a subspace where the laws fail
while the headline reward looks healthy.

\paragraph{Expert-dimension reconstruction.} The internal target
validates the mechanism; the external Benoit expert means test whether
it survives contact with a construct the codes do not determine
(Section~\ref{sec:qsentence-labels}). Training the neural-operator
family at $\ell = 16$ with LLM-generated span supervision at internal
nodes, then scoring roots against the expert means (three seeds per
dimension, test split), gives $r = 0.715 \pm 0.024$ on social,
$0.661 \pm 0.045$ on economic, and $0.642 \pm 0.069$ on environment;
immigration reaches $0.438 \pm 0.094$, and decentralization ($0.244$)
and EU ($0.232$) sit at the floor the CMP-code ceiling predicts for
them. The full-document frontier-LLM read of
Section~\ref{sec:manifesto-llm} sits near $0.86$ on the same task, so
the trained small-substrate stack recovers most of the dense-dimension
signal at a tiny fraction of the per-call cost, and its failure
ordering matches the construct-ceiling ordering rather than being
random. Appendix~\ref{app:qsentence-leaf-ladder} reports the grid and
the winning hyperparameters.

\subsection{Seeding the Summarizer with Gold Quasi-Sentence States}\label{sec:gold-seeding}

The gold labels support one more design worth isolating: using sampled
gold quasi-sentence states to \emph{seed} the summarizer. The question
is whether a learned $g$ can produce a full-document compact state from
a small sample of leaf-window states, so that a fixed learned $f$
matches the document gold: a budgeted version of the local-to-global
reconstruction that Section~\ref{sec:qsentence-labels} showed is
well-posed.

The design has a failure mode we hit and document because it
generalizes. The first implementation fed gold CMP means into the
summarizer's prompt as the sampled state. The resulting $g$ scored
well while learning nothing: it was regressing gold onto gold, and its
outputs degraded when the gold crutch was removed. The corrected
protocol feeds the learned $f$'s own per-leaf predictions
(\texttt{--sample-state-source f\_states}) as the sampled state, so
the pipeline is honest end-to-end: no gold enters at inference time.
Under that protocol, at $\ell = 16$ with $16$ sampled leaf windows per
document, the seeded $g$ reconstructs the full teacher rollup at
all-dimension Pearson $0.953$ ($0.955$ on RILE), against a
sample-only baseline of $0.966$ ($0.949$ on RILE), and agrees with
that baseline at $0.984$: composing learned states recovers
essentially everything the raw sample carries, and loses little to
the additional learned-merge step. Appendix~\ref{app:gold-seeding}
details the protocol, the leak diagnosis, and the per-dimension grid,
and flags the two follow-ups still in flight (the scheduled-sampling
exposure-bias arm and the cross-runtime transfer of gold-trained
prompt programs).

\subsection{What the Quasi-Sentence Experiments Add}

Section~\ref{sec:manifesto-llm} showed that a prompt-only C-Tree
preserves an expert measurement target and leaves an auditable trail.
The quasi-sentence experiments close the loop the controlled sections
promised: on a corpus with genuine node-level gold labels, the same
operator-plus-projection machinery trains, and the places it fails are
the places the label system says it must (sparse dimensions, constructs
the codes do not determine, objectives that do not price the local
laws). The gold quasi-sentence layer is what makes all of this
measurable; corpora without such a layer get the audit of
Section~\ref{sec:estimation} instead, which is the same design run
with sampled oracle calls in place of exhaustive ones.
